\chapter*{General conclusion and perspectives}
\addcontentsline{toc}{chapter}{General conclusion and perspectives}

This final-year project addressed a practical problem faced by many organisations:
internal knowledge exists, but it is hard to find, hard to govern, and hard to use
confidently in daily work. Conducted at \CompanyName{}, the project produced the
Knowledge Platform---a system that brings together document management, controlled
access, semantic search, and cited AI-assisted answers within a single coherent
experience.

The platform delivers a secure multi-role environment in which documents can be
uploaded, indexed, searched, and discussed through persistent conversations.
Retrieval combines semantic and keyword methods so that answers remain anchored
in organisational sources. Administrative and managerial capabilities support user
provisioning, department organisation, notifications, and activity review. From a
validation standpoint, automated tests confirmed core behaviour across authentication,
documents, search, conversations, and governance, while manual scenarios verified
important user-facing flows such as sign-in, document readiness, hybrid Ask, and
profile management.

Relative to the objectives stated in the introduction, the project achieved its
main goals: a role-aware secure platform, improved retrieval over heterogeneous
documents, intelligent assistance with source transparency, structured validation,
and a reproducible deployment model. Remaining limitations include reliance on an
external language-model provider, the absence of large-scale performance
benchmarking, and evaluation sensitivity to corpus coverage and quality.

Beyond technical delivery, the project strengthened competencies in requirements
engineering, system design, information retrieval, AI integration, security
awareness, and critical evaluation of intelligent systems---skills directly
relevant to enterprise software engineering and data-oriented practice.

Future work could extend the platform through cloud object storage, enterprise
single sign-on, on-premise model options, tighter continuous-evaluation pipelines
for answer quality, mobile access, and field deployment with real corpus
onboarding and formal user acceptance testing at the host organisation.
